% ECONOMICS_PROBLEM: {"id": "J44-623", "title": "Occupational licensing", "jel_code": "J44", "source_id": "OP-0025"}
% !TeX program = xelatex
\documentclass[11pt,a4paper]{article}
\usepackage[margin=23mm]{geometry}
\usepackage{fontspec}
\setmainfont{Latin Modern Roman}
\usepackage{amsmath,amssymb,mathtools}
\usepackage{xcolor,enumitem,booktabs,longtable,array,xurl,hyperref}
\definecolor{muted}{HTML}{53636C}
\hypersetup{colorlinks=true,urlcolor=blue,linkcolor=blue}
\setlength{\parindent}{0pt}
\setlength{\parskip}{5pt}
\newcommand{\R}{\mathbb R}
\newcommand{\E}{\mathbb E}
\newcommand{\N}{\mathbb N}
\newcommand{\ind}{\mathbf 1}
\newcommand{\Var}{\operatorname{Var}}
\newcommand{\Cov}{\operatorname{Cov}}
\newcommand{\supp}{\operatorname{supp}}
\DeclareMathOperator*{\argmin}{arg\,min}
\DeclareMathOperator*{\argmax}{arg\,max}
\newcommand{\entrymeta}[3]{{\sffamily\small\color{muted}\textbf{\texttt{#1}}\quad|\quad #2\par #3\par}}
\newcommand{\Problemref}[1]{\texttt{#1}}
\newcommand{\JELref}[2]{\texttt{#2} (JEL \texttt{#1})}
\begin{document}
\section*{Occupational licensing}
\textbf{Problem ID:} \texttt{J44-623}\quad \textbf{JEL:} \texttt{J44}\par
% BEGIN ECONOMICS STATEMENT
\label{problem:OP-0025}
\label{atlas:prob:L5}
\noindent\textbf{Original atlas ID:} L5 (entry 25). \textbf{Classification:} Editorial equilibrium welfare and policy-design problem.

\noindent\textbf{Original atlas question:} When does licensing improve quality enough to justify entry and mobility restrictions?

\subsection*{Definitions and assumptions}
Fix an occupation, jurisdiction, and policy horizon. Let $\ell\in\mathcal L$ specify licensing requirements, enforcement, reciprocity, and grandfathering. A market model $M$ maps $\ell$ into provider entry, effort, service quality $q$, prices, consumer choice, and external harms. Quality may be a vector containing independently defined safety, effectiveness, and timeliness outcomes. Specify information asymmetries, demand, production costs, and government costs. Assume equilibrium existence and either uniqueness or a selection rule. Let $B_i(q)$ be consumer $i$'s monetized service benefit, $H_M(\ell)$ external harms, $C_M(\ell)$ real production and entry costs, and $A_M(\ell)$ administration costs, all measured on a common discounted basis.
\subsection*{Formal research question}
With fixed population and social weights, define total surplus
\[
 W_M(\ell)=\mathbb E\left[\sum_i B_i(q_i(\ell))\right]
             -H_M(\ell)-C_M(\ell)-A_M(\ell).
\]
Prices paid to domestic providers are transfers in this aggregate unless additional distributional weights, foreign ownership, or fiscal distortions are specified. Identify or sharply bound $W_M(\ell)-W_M(\ell_0)$ for each feasible reform $\ell$, allowing provider composition to change. Conditional on announced model and welfare assumptions, characterize $\arg\max_{\ell\in\mathcal L}W_M(\ell)$ when the maximum exists; otherwise characterize the supremum and approximately optimal policies. Determine which measured quality improvements justify their resource and exclusion costs.
\subsection*{Interpretation and scope}
A licensed-unlicensed wage comparison does not identify consumer benefits, entry effects, or net surplus. Provider-level quality conditional on survival is also insufficient when licensing changes which providers operate and which consumers receive services. Border or reform designs require treatment consistency, credible comparison regions, and control or modeling of cross-border service substitution. An instrumental-variable effect for affected providers does not automatically identify an equilibrium reform. Safety constraints can be imposed separately as $H_M(\ell)\leq\bar H$ for a specified threshold $\bar H$; that normative constraint is not inferred from earnings data. Distributional analysis should report access and prices for consumer groups, entry opportunities, and provider incomes in addition to aggregate surplus.
\subsection*{Source audit and status}
All three cited records were located. [S2] is the 2019 working paper, with a 2023 journal publication, and already undertakes a welfare analysis. Therefore welfare evaluation itself is not presented as wholly unstudied. The editorial extension requires multidimensional quality and changing provider composition across specified reforms. No general open-status certification follows from the dates of [S1]--[S3].

\subsection*{References}
\begin{enumerate}[label=\textnormal{[S\arabic*]},leftmargin=*]
\item Morris M. Kleiner and Alan B. Krueger (2013). \emph{Analyzing the Extent and Influence of Occupational Licensing on the Labor Market}. Journal of Labor Economics 31(S1), S173--S202. \url{https://doi.org/10.1086/669060}. \textit{Locator:} Primary publisher, working-paper, or author record: bibliographic information and abstract; full-text theorem verification is not claimed. \textit{Role:} Licensing prevalence and wage evidence; wage associations alone are not welfare effects.
\item Morris M. Kleiner and Evan J. Soltas (2019). \emph{A Welfare Analysis of Occupational Licensing in U.S. States}. NBER Working Paper 26383; published in The Review of Economic Studies 90(5), 2481--2516 (2023). \url{https://www.nber.org/papers/w26383}. \textit{Locator:} Primary publisher, working-paper, or author record: bibliographic information and abstract; full-text theorem verification is not claimed. \textit{Role:} Structural welfare analysis of occupational licensing. The atlas year refers to the working-paper version; the 2023 journal version is verified in turn5search27.
\item Peter Q. Blair and Bobby W. Chung (2019). \emph{How Much of Barrier to Entry is Occupational Licensing?}. British Journal of Industrial Relations 57(4), 919--943. \url{https://doi.org/10.1111/bjir.12470}. \textit{Locator:} Primary publisher, working-paper, or author record: bibliographic information and abstract; full-text theorem verification is not claimed. \textit{Role:} Boundary-discontinuity evidence on occupational choice and entry.
\end{enumerate}

\subsection*{Common conventions: Source mathematical conventions}

Unless an entry states otherwise, agent and firm sets are finite. State and action spaces are measurable, strategies and outcome maps are measurable, probability laws are specified, and expectations exist for the targets under discussion. Infinite-horizon discount factors lie in $(0,1)$ and discounted objectives are finite. Norms, tolerances and complexity input sizes must be specified for computational comparisons. Constants or primitives that appear locally are inputs, not empirical facts asserted by this catalogue.

\textbf{Identification.} A statistical model $m\in\mathcal M$ implies an observable law $P_m$ and target $\theta(m)$. At law $P$, its identified set is
\[
\Theta(P;\mathcal M)=\{\theta(m):m\in\mathcal M,\ P_m=P\}.
\]
Point identification holds when this set is a singleton, or equivalently when $P_{m_1}=P_{m_2}$ implies $\theta(m_1)=\theta(m_2)$ for all compatible models. Sharp bounds mean bounds attained, or approached when the boundary is not attained, by compatible models. Writing this set is a definition, not a solution: the substantive task is to establish informative restrictions and derive or estimate the set. An unrestricted-confounding model may give only trivial bounds.

\textbf{Inference.} Identification, estimability and valid uncertainty quantification are separate questions. Fix $\alpha\in(0,1)$. A confidence procedure $C_n$ over a declared law class $\mathcal P$ has asymptotic uniform coverage at level $1-\alpha$ when
\[
\liminf_{n\to\infty}\inf_{P\in\mathcal P}\Pr_P\{\theta(P)\in C_n\}\geq1-\alpha.
\]
For set-identified targets the coverage event must be defined separately, for example coverage of the full identified set. If an entry uses identification-indexed models instead of uniquely indexed laws, the same coverage convention is applied over those models. Pointwise limits do not establish uniform validity, and asymptotic validity does not guarantee finite-sample precision. Sample sizes, dependence structures and weak-identification sequences are part of each application.

\textbf{Counterfactuals.} $Y_i(a)$ is a potential outcome under a specified intervention, including its timing, implementation and versions. It is not $Y_i$ conditional on voluntarily choosing $a$. A full intervention vector permits interference; shorthand scalar treatment notation does not silently impose no interference. Assignment, selection, attrition, anticipation and measurement must be specified. A claim about an instrument's local effect is not automatically a population policy effect.

\textbf{Economic equilibrium.} A counterfactual model includes best responses, constraints, feasibility, market clearing, state transitions and expectations consistent with the stated information. When several equilibria exist, either a declared selector is an input or the target is set-valued. Comparisons across policy regimes must use the stated selection convention. Reduced-form relationships are not presumed invariant after an equilibrium-changing intervention.

\textbf{Optimization and welfare.} Optimization is over the stated feasible set. If attainment is not established, $\sup$ or $\inf$ and $\varepsilon$-optimal choices replace an asserted maximizer or minimizer. For example, with value $V=\sup_{a\in\mathcal A}W(a)<\infty$, an $\varepsilon$-optimal set is $\{a\in\mathcal A:W(a)\geq V-\varepsilon\}$ for $\varepsilon>0$. Benefits, costs, risk and health or environmental outcomes can be aggregated only using declared units and conversion weights. Interpersonal utility weights, the treatment of fiscal transfers and foreign profits, discounting, and the behavioral welfare criterion are explicit inputs. A comparative-static derivative additionally requires existence and appropriate differentiability of the selected counterfactual.

\textbf{Sources.} Local labels [S1], [S2], and so forth restart in every question. Locators identify the relevant abstract, model, section or result at the level actually checked; a bibliographic or abstract-level verification is not represented as inspection of an entire proof. Working papers are distinguished from later publications and changing versions. Source summaries are paraphrased. All URL checks and editorial formulations refer to the audit date stated above.
% END ECONOMICS STATEMENT
\end{document}
